\documentclass{article}
\usepackage{iclr2026_conference,times}
\input{math_commands.tex}
\usepackage{hyperref}
\usepackage{url}
\usepackage{booktabs}
\title{Planner Recoverability Scores: ICLR Main-Conference Gate Archive}
\author{Anonymous Authors}
\begin{document}
\maketitle
\begin{abstract}
This document is not an ICLR main-conference submission. It is the v3 main-conference gate for a generated robotics draft about task and motion planning. The original research bet was: Score plans by recoverability under physical deviation, not just nominal feasibility. After a stricter hostile review, the honest terminal decision is \textbf{KILL/ARCHIVE}. Submission-hardening version: v3.
\end{abstract}

\section{Gate Decision}
The v2 paper was workshop-only. Under the active ICLR-main standard, that is insufficient. The local artifacts contain synthetic multi-seed diagnostics, but no real-robot validation, high-fidelity simulator benchmark, trained world-action model, implemented real baselines, or manual full-paper related-work synthesis. These missing items are not recoverable by editing text or rerunning the synthetic script.

\section{Fatal ICLR-Main Blockers}
\begin{itemize}
\item Evidence is synthetic and generated from a shared batch template.
\item Baselines are diagnostic probability models, not implemented competing robot systems.
\item Ablations do not correspond to a trained model checkpoint.
\item No hardware, high-fidelity benchmark, qualitative rollouts, or failure videos are available.
\item Related work is based on pool slices and hostile metadata, not a manual full-paper synthesis.
\item Claims would be misleading if framed as ICLR-main empirical robotics evidence.
\end{itemize}

\section{Hostile Prior-Work Pressure}
The closest local hostile records include:
\begin{itemize}
\item XFlowMP: Task-Conditioned Motion Fields for Generative Robot Planning with Schrodinger Bridges (2025)
\item An incremental constraint-based framework for task and motion planning (2018)
\item Combined Task and Motion Planning via Sketch Decompositions (2024)
\item Representation, learning, and planning algorithms for geometric task and motion planning (2022)
\item Socially intelligent task and motion planning for human-robot interaction (2020)
\end{itemize}
These works make generic world models, uncertainty, model-based planning, and robot policy robustness crowded. Without real empirical evidence, the remaining novelty boundary is too weak for ICLR main.

\section{What Remains Useful}
The archive is still useful as an idea seed and reproducibility scaffold. The synthetic code, stress-test CSVs, reviewer attacks, and novelty-boundary docs can inform a future project. The correct next step is not more paper polishing; it is new evidence: implemented models, real baselines, manual related work, and robot or high-fidelity simulator validation.

\section{Terminal Action}
Decision: \textbf{KILL/ARCHIVE for ICLR main}. Do not submit this paper to ICLR main in the current form. Revival requires a substantially new empirical project.

\end{document}
